\section{Introduction}
\label{f:sec:introduction}

Fix one viscosity \(\nu>0\), one divergence-free datum
\(u_0\in H^\infty(\R^3;\R^3)\), and one prescribed source
\[
 f\in C^\infty([0,\infty);\mathcal S(\R^3;\R^3)).
\]
Forced local theory produces a unique pressure-normalized classical history
\((u,p)\) on a maximal interval \([0,T_*)\). Part~\ref{paper:part-forced} follows that one
history. Its question is whether a finite value of \(T_*\) can terminate the
solution generated by \((u_0,\nu,f)\).

\begin{maintheorem}
\label{f:thm:main-global}
For every \(\nu>0\), every \(u_0\in\HsSigma(\R^3)\), and every
\(f\in C^\infty([0,\infty);\mathcal S(\R^3;\R^3))\), there is a unique
pressure-normalized pair
\[
 (u,p)\in C^\infty([0,\infty)\times\R^3;\R^3\times\R)
\]
satisfying
\[
 \partial_tu+(u\cdot\nabla)u+\nabla p=\nu\Delta u+f,
 \qquad \nabla\cdot u=0,
 \qquad u(0)=u_0,
 \qquad \lim_{\abs{x}\to\infty}p(t,x)=0.
\]
It obeys the forced energy identity
\[
 \frac12\norm{u(t)}_{L^2}^2
 +\nu\int_0^t\norm{\nabla u(s)}_{L^2}^2\dd s
 =\frac12\norm{u_0}_{L^2}^2
 +\int_0^t\ip{f(s)}{u(s)}_{L^2}\dd s
 \qquad(t\ge0).
\]
\end{maintheorem}

\begin{corollary}[Weak--strong uniqueness and pressure normalization]
\label{f:cor:main-weak-strong}
For the solution of the preceding theorem, every forced Leray--Hopf pair
\((v,q)\) with that datum, viscosity, and
prescribed source has \(v=u\). Distributionally,
\(q=p+c(t)\) for a scalar distribution \(c\); if
\(q(t,\cdot)\in C_0(\R^3)\) almost everywhere, then \(q=p\) almost
everywhere.
\end{corollary}

Leray's weak solutions establish global finite-energy existence
\cite{Leray1934}. The regularity results of Serrin
\cite{Serrin1962}, Caffarelli--Kohn--Nirenberg \cite{CKN1982}, and
Escauriaza--Seregin--\v{S}ver\'ak \cite{ESS2003} identify increasingly sharp
ways in which a singularity would have to evade classical control.
Fujita--Kato and Kato provide the local strong theory used both at the
initial time and at the terminal restart \cite{FujitaKato1964,Kato1984}.
The proof below works directly with the complete derivative structure of the
prescribed forced history.

Forcing changes every level of that structure. Its solenoidal part
\(g=\Pp f\) enters the projected velocity equation. Its gradient part is
absorbed by the normalized pressure,
\[
 -\Delta p=\partial_i\partial_j(u_i u_j)-\nabla\cdot f,
 \qquad
 p=R_iR_j(u_i u_j)-(-\Delta)^{-1}\nabla\cdot f.
\]
Every temporal row contains the corresponding derivative
\(F_n=\partial_t^nf\); every energy row contains force work; and a restart at
\(T_*\) retains that absolute-time source. These are parts of a single
forced velocity--pressure--source history.

The prescribed datum and complete equation generate the forced mixed jet
with its spatial Sobolev intersection. Compatible temporal rows give the
mixed intersection on the original interval. Critical height reads the
spatial domains during generation; completed-height substitution then
displays the viscous relationships with transport, pressure, and force work
retained. Finite rectangles select the required coordinates and relationships
from this constructed jet. Their adjacent rows reach one smooth endpoint.
The forced recurrence fixes its complete pressure and velocity jet, and local
theory restarts from \(u_*=u(T_*)\) with the prescribed absolute-time force.
Matching the endpoint jets extends the incoming history and contradicts a
finite maximal time.

Section~\ref{f:sec:forced-argument-overview} derives this architecture in one
view. The analytic framework and the mixed-jet construction follow, then finite
coordinate selections and their mixed-index exhaustion. The independently
derived critical-height and viscous identities give relationships among the
constructed jet's coordinates. Both descriptions retain the intersection used
in the endpoint continuation.
The remaining sections record matrix formulations and quantitative
consequences of these finite rectangles.
